%! Multiplying \& Dividing Rational Expressions — Unit 4, Lesson 4.2 — Homework
\documentclass[10pt]{article}
\usepackage{algebra2-article}
\usepackage{algebra2-boxes}
\newcommand{\TallMath}[1]{$\displaystyle #1\rule[-1.4em]{0pt}{3.2em}$}

\begin{document}

\pageheader{Unit 4, Lesson 4.2}{Homework}
\namedateperiod

\vspace{0.10in}

\begin{notesbox}{Practice}
Every answer needs \emph{two} parts: the simplest form \textbf{and} the restrictions. Check all
\textbf{three} sources on every division --- both denominators \emph{and} the divisor's numerator.

\begin{enumerate}[leftmargin=1.9em, itemsep=11pt]
  \item \textbf{Monomial factors.} Simplify and restrict.
    \begin{enumerate}[label=(\alph*), leftmargin=2.1em, itemsep=5pt]
      \item $\dfrac{4x^3}{9y}\cdot\dfrac{3y^4}{8x}=\;$\blank{2.2cm} \quad restrictions: \blank{2.6cm}
      \item $\dfrac{5a^2}{6b^3}\div\dfrac{10a}{3b}=\;$\blank{2.2cm} \quad restrictions: \blank{2.6cm}
      \item $\dfrac{7x^2y}{2}\div\dfrac{14xy^3}{5}=\;$\blank{2.2cm} \quad restrictions: \blank{2.6cm}
    \end{enumerate}
    \vspace{2pt}
    In part (c), neither original denominator can ever be zero. So where did the restrictions come
    from? \blank{5.6cm}

  \item \textbf{Binomial and quadratic factors.} Factor completely first. One of these hides a pair of
        \emph{opposite} binomials --- watch the sign.
    \begin{enumerate}[label=(\alph*), leftmargin=2.1em, itemsep=7pt]
      \item $\dfrac{x+3}{x-2}\cdot\dfrac{x^2-4}{x^2-9}=\;$\blank{2.4cm} \quad $x\neq$ \blank{1.8cm}
      \item $\dfrac{x^2+6x+8}{x^2-1}\cdot\dfrac{x-1}{x+4}=\;$\blank{2.4cm} \quad $x\neq$ \blank{1.8cm}
      \item $\dfrac{x^2-36}{5x}\div\dfrac{x+6}{10x^2}=\;$\blank{2.4cm} \quad $x\neq$ \blank{1.8cm}
      \item $\dfrac{x^2+2x-15}{x^2+x-6}\div\dfrac{x+5}{x+3}=\;$\blank{2.4cm} \quad $x\neq$ \blank{2.2cm}
      \item $\dfrac{x^2-49}{x+1}\cdot\dfrac{3x+3}{7-x}=\;$\blank{2.4cm} \quad $x\neq$ \blank{1.8cm}
      \item $\dfrac{x^2-2x}{x^2+5x+6}\div\dfrac{x^2-4}{x+3}=\;$\blank{2.4cm} \quad $x\neq$ \blank{2.2cm}
    \end{enumerate}

  \item \textbf{A rectangle.} A rectangle has width $\dfrac{2x}{x+5}$ units and length
        $\dfrac{x^2+5x}{4}$ units.
    \begin{enumerate}[label=(\alph*), leftmargin=2.1em, itemsep=5pt]
      \item Area $=$ width $\times$ length $=$ \blank{2.4cm}, \; with restriction $x\neq$ \blank{1.2cm}
      \item Your area has no denominator containing $x$. Explain why the restriction still belongs to
            it. \blank{5.0cm}
      \item Both dimensions must be positive for a real rectangle. Which values of $x$ make sense?
            Write your answer as a union of intervals. \blank{4.2cm}
    \end{enumerate}

  \item \textbf{Justify.} Are $\dfrac{x+1}{x-2}\div\dfrac{x+1}{x-3}$ and $\dfrac{x-3}{x-2}$ the same
        expression? Answer in terms of \emph{inputs}: where do they agree, where do they disagree, and
        what does each graph look like?
        \writelines{3}
\end{enumerate}
\end{notesbox}

\begin{extensionbox}
\textbf{Part 1 --- Three students, three answers.} Each student simplified
$\dfrac{x^2-1}{x+4}\div\dfrac{x-1}{x^2-16}$ and got something different. Decide who is right, and name
each error.
\begin{center}
\renewcommand{\arraystretch}{1.4}
\begin{tabularx}{\linewidth}{@{}l l X@{}}
\hline
\textbf{Student} & \textbf{Answer} & \textbf{Right or wrong --- and why?} \\
\hline
Dev   & $(x+1)(x-4)$, $x\neq-4,4,1$ & \blank{5.4cm} \\
Elin  & $(x+1)(x-4)$, no restrictions & \blank{5.4cm} \\
Farid & $\dfrac{(x-1)^2(x+1)}{(x+4)^2(x-4)}$ & \blank{5.4cm} \\
\hline
\end{tabularx}
\end{center}

\vspace{4pt}
\textbf{Part 2 --- Design one.} Write a \emph{division} problem whose simplified answer is $x+2$ and
whose restrictions are exactly $x\neq3$, $x\neq0$, and $x\neq-1$.
\begin{center}
\blank{3.6cm} $\div$ \blank{3.6cm}
\end{center}
\begin{enumerate}[label=(\alph*), leftmargin=1.9em, itemsep=5pt]
  \item Which factor did you put in the divisor's \emph{numerator}, and which restriction did that buy
        you?
  \item Your answer $x+2$ is a polynomial --- its graph is a line. What does that line look like once
        the restrictions are applied?
\end{enumerate}
\writelines{3}
\end{extensionbox}

\begin{spiralbox}
\textbf{Coming next --- Lesson 4.3: Adding \& Subtracting Rational Expressions (and complex
fractions).} Multiplying and dividing never needed a common denominator; adding and subtracting need
nothing else. You will factor every denominator (same first step as always), build the \textbf{least
common denominator} out of those factors, and rewrite each fraction to match it. Then the payoff:
a \textbf{complex fraction} --- a fraction whose numerator and denominator are themselves fractions ---
is nothing more than \emph{combine the top, combine the bottom, then divide}. And dividing is exactly
what you did today, three restriction sources and all.
\end{spiralbox}

\end{document}
